\scanpage{22}
% archive foliation: 21
\[
\operatorname{Hom}_k(\Lambda^i\Psi_*,\Lambda^j\Psi_*)\simeq
\left\{\begin{array}{ll}0&\text{si }j\neq i\\k&\text{si }j=i\end{array}\right.
\]
de façon précise
\[
\operatorname{Hom}_k(\Lambda^i\Psi_*,\Lambda^i\Psi_*)
\simeq k\operatorname{id}_{\Lambda^i\Psi_*}\simeq k,
\]
et
\[
\operatorname{Hom}_k(\Lambda^i\Psi_*,\Phi_{\deleted{j}}(j,n))\simeq
\left\{\begin{array}{ll}
0&\text{si }j\neq\deleted{i,}\ i+1\\
k&\text{si }j=i+1
\end{array}\right.
\]
en particulier
\[
\operatorname{Hom}_k(\Lambda^i\Psi_*,\Lambda^j\Phi_*)\simeq
\left\{\begin{array}{ll}0&\text{si }j\neq i+1\\k&\text{si }j=i+1\end{array}\right.
\]
et
\[
\operatorname{Hom}_k(\Lambda^i\Psi_*,\Lambda^{i+1}\Phi_*)
=k\cdot\iota\simeq k
\]
où $\iota:\Lambda^i\Psi_*\hookrightarrow\Lambda^{i+1}\Phi_*$ est l'inclusion
canonique.

\uline{NB} On a un \inserted{hom de} suites exactes can
\diagram{\begin{tikzcd}[column sep=large]
0\arrow[r,dash]&\Lambda^{i-1}\Psi_*\arrow[r,dash]\arrow[d,"n\operatorname{id}"']&
\Lambda^{i\deleted{\unsure{+1}}}\Phi_*\arrow[r]\arrow[d,equal]&
\Lambda^i\Psi_*\arrow[r,dash]\arrow[d,Rightarrow,"\operatorname{id}"]&0\\
0\arrow[r,dash]&\Lambda^{i-1}\Psi_*\arrow[r]&
\deleted{\unsure{\Lambda^i}}\Phi_*(i,n)\arrow[r,dash]&\Lambda^i\Psi_*\arrow[r,dash]&0
\end{tikzcd}}
qui montre que $\Phi_*(i,n)$ est l'ext.
déduite de \inserted{l'ext} $\Lambda^i\Phi_*$ de $\Lambda^i\Psi_*$ par $\Lambda^{i-1}\Psi_*$
par l'hom $n\operatorname{id}:\Lambda^{i-1}\Psi_*\;-\;\Lambda^{i-1}\Psi_*$.
